\uuid{N0bE}
\exo7id{1685}
\titre{exo7 1685}
\auteur{ortiz}
\organisation{exo7}
\datecreate{1999-04-01}
\isIndication{false}
\isCorrection{false}
\chapitre{Réduction d'endomorphisme, polynôme annulateur}
\sousChapitre{Sous-espace stable}
\cpge{mpsi-mp2i.dimension-finie.b, pcsi.dimension-finie.b, ptsi.dimension-finie.b, mpsi-mp2i.dimension-finie.c, pcsi.dimension-finie.c, ptsi.dimension-finie.c, mp-mpi.reduction-endomorphismes.b, pc.reduction.a, psi.reduction-endomorphismes.a, pt.reduction-endomorphismes.a, pc.complements-espaces-vectoriels.b, psi.complements-espaces-vectoriels.b, pt.complements-algebre-lineaire.b, bcpst2.valeurs-propres-vecteurs-propres.a}
\module{Algèbre}
\niveau{L2}
\difficulte{}

\contenu{
\question{
Soient $E$ un $\Rr$-espace vectoriel de dimension
finie, $f$ un endomorphisme de $E$. Montrer que si
$f$ admet un sous-espace vectoriel propre de
dimension $p\geq 2$ alors il admet une infinit\'e
de sous-espaces vectoriels stables par $f$.
}
}
